\chapter{窗口型无零点定理}
\label{chap:window-proof}
\label{chap:window}

本章证明一个明确的解析推断：如果归一化的平滑素数和在任意远处的整段窗口上受到控制，那么相应区域内就不能有 $\zeta$ 零点。窗口估计是这一推断的输入；下面完整证明从该输入到无零点结论的过程。

沿用前文的实值函数 $G$、$g(u)=G(e^u)$ 及其拉普拉斯变换 $\Phi$。本章使用前文已经建立的以下性质：$g$ 连续，存在固定 $C_0>0$ 使
\begin{equation}
 |g(u)|\le C_0e^{2u}\qquad(u\ge0);
 \label{eq:win-growth}
\end{equation}
当 $\operatorname{Re}s>2$ 时，
\begin{equation}
 \Phi(s)=\int_0^\infty g(u)e^{-su}\,du.
 \label{eq:win-laplace}
\end{equation}
此外，$\Phi$ 亚纯延拓到半平面
\[
 \mathcal H=\{s\in\mathbb C:\operatorname{Re}s>0\}.
\]
它的极点集合 $Z$ 恰为 $\zeta$ 在 $\mathcal H$ 中的零点集合。每个 $\rho\in Z$ 都是 $\Phi$ 的单极点，留数 $c_\rho\ne0$；$\Phi$ 在正实轴上没有极点。这里 $\zeta$ 的 $s=1$ 极点已经在构造 $\Phi$ 时被消去。

\begin{theorem}[窗口型无零点定理]
\label{thm:window-zeta}
固定 $0<\theta<1$。假定
\begin{equation}
 \forall Y_0\ge1\quad
 \exists Y\ge Y_0\quad
 \forall x\in[Y,Y^{3/2}],
 \qquad G(x)\le x^\theta.
 \label{eq:win-pointwise}
\end{equation}
则 $\Phi$ 在 $\operatorname{Re}s>\theta$ 内没有极点。因此，
\[
 \operatorname{Re}s>\theta,\qquad s\ne1
 \quad\Longrightarrow\quad \zeta(s)\ne0.
\]
\end{theorem}

证明将比较同一组高阶导数的两个界。窗口估计给出上界；若存在一个待排除的极点，它又会迫使这些导数具有更大的下界。关键是把两种估计放在同一尺度上。

\section{把窗口上界变成导数上界}

令 $U=\log Y$。幂区间 $[Y,Y^{3/2}]$ 变成对数区间
\[
 W_U=[U,3U/2].
\]
式 \eqref{eq:win-pointwise} 意味着：对任意 $U_0\ge1$，总能选到 $U\ge U_0$，使
\begin{equation}
 g(u)\le e^{\theta u}\qquad(u\in W_U).
 \label{eq:win-log-window}
\end{equation}
这是因为可先对起点 $Y_0=e^{U_0}$ 使用窗口假设，再取对数。以后称满足 \eqref{eq:win-log-window} 的 $U$ 为好尺度。所需的只是好尺度无界。

对整数 $k\ge0$ 和 $\operatorname{Re}z>2$，定义
\begin{equation}
 A_k(z)=\frac{(-1)^k}{k!}\Phi^{(k)}(z)
       =\frac1{k!}\int_0^\infty g(u)u^ke^{-zu}\,du.
 \label{eq:win-Ak}
\end{equation}
第二个等式可以逐次在积分号下求导。确实，在任意这样的 $z$ 的一个小邻域内，实部都有统一下界 $\sigma_0>2$，而第 $k$ 阶被积函数的绝对值由
\[
 C_0u^ke^{-(\sigma_0-2)u}
\]
控制。它在半轴上可积，因而每一阶求导都有所需的可积控制函数。这也说明 $\Phi$ 在 $\operatorname{Re}z>2$ 全纯。

固定
\begin{equation}
 \sigma=336.
 \label{eq:win-sigma}
\end{equation}
所有好尺度都使用这个相同的实部。$336$ 只是让后面的两侧尾估计同时成立的一种方便选择；下面会直接验证这些数值，不把它们当作额外假设。对 $u\ge0$，记
\[
 w_k(u)=\frac{u^ke^{-\sigma u}}{k!}\ge0.
\]
由于 $g$ 实值，$A_k(\sigma)$ 也是实数。三角不等式给出
\begin{align}
 |A_k(\sigma+it)|+\operatorname{Re}A_k(\sigma)
 &\le \int_0^\infty \bigl(|g(u)|+g(u)\bigr)w_k(u)\,du.
 \label{eq:win-positive-part}
\end{align}
这一步是使用单侧上界的关键。对于实数 $v$，
\[
 |v|+v=2\max\{v,0\}.
\]
因此只要控制 $g$ 的上方，就能控制 \eqref{eq:win-positive-part} 的右边。若把左边擅自改成两个绝对值之和，就失去了这项性质。

对任意 $a>0$，指数函数的幂级数还给出
\begin{equation}
 \frac{u^ke^{-au}}{k!}\le a^{-k}\qquad(u\ge0),
 \label{eq:win-kernel-max}
\end{equation}
因为 $e^{au}\ge(au)^k/k!$。在好窗口内，由
$|g|+g\le2e^{\theta u}$ 和 \eqref{eq:win-kernel-max}，得到
\[
 \int_{W_U}(|g|+g)w_k
 \le \frac{2}{(\sigma-\theta)^k}\,\frac U2
 =\frac{U}{(\sigma-\theta)^k}.
\]
窗口外则用全局粗界 \eqref{eq:win-growth}。于是对每个好尺度 $U$、每个 $k\ge0$ 及每个 $t\in\mathbb R$，
\begin{equation}
 |A_k(\sigma+it)|+\operatorname{Re}A_k(\sigma)
 \le \frac{U}{(\sigma-\theta)^k}+2C_0T_{U,k},
 \label{eq:win-before-tail}
\end{equation}
其中
\[
 T_{U,k}=
 \int_{[0,\infty)\setminus W_U}
       \frac{u^ke^{-334u}}{k!}\,du.
\]
接下来的工作是控制这一尾积分。我们不能因为窗口外没有算术上界，就把那部分积分删去。

\section{用明确参数把积分核集中到窗口内}

对每个 $U\ge1$，取
\begin{equation}
 k=\lceil420U\rceil.
 \label{eq:win-k-choice}
\end{equation}
于是
\[
 420U\le k\le420U+1,\qquad k\ge U.
\]
导数阶数与窗口位置一起增大。相应权核的主要贡献落在窗口内部；下面的计算对窗口外剩余的贡献给出统一界。

先记录一个基本积分公式：
\begin{equation}
 \frac1{k!}\int_0^\infty u^ke^{-au}\,du
       =a^{-(k+1)}\qquad(a>0).
 \label{eq:win-gamma-integral}
\end{equation}
它从 $k=0$ 的指数积分出发，逐次分部积分即可得到；指数衰减使无穷远处的边界项为零。

在窗口右侧 $u\ge3U/2$，
\[
 e^{-334u}=e^{-320u}e^{-14u}
       \le e^{-21U}e^{-320u}.
\]
在窗口左侧 $0\le u\le U$，
\[
 e^{-334u}=e^{-348u}e^{14u}
       \le e^{14U}e^{-348u}.
\]
将两部分非负积分分别扩大到整个半轴，并使用
\eqref{eq:win-gamma-integral}，得到
\begin{equation}
 T_{U,k}\le
 \frac{e^{-21U}}{320^{k+1}}
 +\frac{e^{14U}}{348^{k+1}}.
 \label{eq:win-tail-two}
\end{equation}

为了与导数上界比较，把两项分别乘以 $\sigma^k=336^k$。由
$\log(1+x)\le x$（$x>-1$）和 \eqref{eq:win-k-choice}，
\begin{align*}
 336^k\frac{e^{-21U}}{320^{k+1}}
 &=\frac1{320}
   \exp\left(-21U+k\log\frac{21}{20}\right)\\
 &\le\frac1{320}\exp(-21U+k/20)
 \le\frac{e^{1/20}}{320}
 \le\frac e{320}.
\end{align*}
同样，
\begin{align*}
 336^k\frac{e^{14U}}{348^{k+1}}
 &=\frac1{348}
   \exp\left(14U+k\log\frac{28}{29}\right)\\
 &\le\frac1{348}\exp(14U-k/29)\\
 &\le\frac1{348}\exp(-14U/29)
 \le\frac1{348}.
\end{align*}
因此可以固定
\begin{equation}
 C_1=\frac e{320}+\frac1{348},
 \qquad
 T_{U,k}\le\frac{C_1}{\sigma^k}.
 \label{eq:win-tail-final}
\end{equation}
这里的 $C_1$ 与窗口位置 $U$ 无关。

令 $C_2=2C_0C_1$。由于 $0<\theta<1<\sigma$，
有 $\sigma^{-k}\le(\sigma-\theta)^{-k}$。
将 \eqref{eq:win-tail-final} 代入
\eqref{eq:win-before-tail}，得到本章的统一上界：
\begin{equation}
 \boxed{\;
 |A_k(\sigma+it)|+\operatorname{Re}A_k(\sigma)
 \le\frac{U+C_2}{(\sigma-\theta)^k},
 \qquad k=\lceil420U\rceil.
 \;}
 \label{eq:win-master-upper}
\end{equation}
它对每个好尺度 $U\ge1$ 和所有实数 $t$ 同时成立。

\section{在极点附近选出能够比较的两个圆盘}

现在反设存在 $\rho_0\in Z$，满足
$\operatorname{Re}\rho_0>\theta$。我们将选出一个实中心圆盘，其中没有极点；再竖直移动圆心，得到一个只含一个极点的圆盘。两个圆心都具有相同实部 $\sigma$，因此可以一起代入
\eqref{eq:win-master-upper}。

首先，由 $\Phi$ 在 $\operatorname{Re}s>2$ 全纯可知
\begin{equation}
 \operatorname{Re}\rho\le2<\sigma\qquad(\rho\in Z).
 \label{eq:win-poles-left}
\end{equation}
亚纯函数的极点在 $\mathcal H$ 内没有聚点，所以 $\mathcal H$ 中每个紧集只含有限个极点。

记 $R=\sigma-\theta$。闭圆盘
$\overline B(\sigma,R)$ 的实部下界为 $\theta>0$，故它紧含于 $\mathcal H$，其中的极点构成有限集。

如果该闭圆盘没有极点，取 $D=R$、$\rho_1=\rho_0$。此时 $B(\sigma,D)$ 无极点，且
\[
 0<\sigma-\operatorname{Re}\rho_1<D.
\]
如果闭圆盘中有极点，就在这有限多个极点中取一个离 $\sigma$ 最近的 $\rho_1$，并令
\[
 D=|\sigma-\rho_1|\le R.
\]
因为 $\sigma$ 不是极点，所以 $D>0$；最小性的定义保证开圆盘 $B(\sigma,D)$ 无极点。又因为正实轴没有极点，$\rho_1$ 非实，因而
\[
 D^2=(\sigma-\operatorname{Re}\rho_1)^2
             +(\operatorname{Im}\rho_1)^2
       >(\sigma-\operatorname{Re}\rho_1)^2.
\]
结合 \eqref{eq:win-poles-left}，仍得
$0<\sigma-\operatorname{Re}\rho_1<D$。

两种情形都给出了
\begin{equation}
 \begin{gathered}
 0<D\le\sigma-\theta,\qquad
 B(\sigma,D)\cap Z=\varnothing,\\
 \rho_1\in Z,\qquad
 0<a:=\sigma-\operatorname{Re}\rho_1<D.
 \end{gathered}
 \label{eq:win-free-disc}
\end{equation}
这一构造只需要检查紧圆盘中的有限个极点。

\subsection{竖移圆心并打破等距}

令
\[
 z_0=\sigma+i\operatorname{Im}\rho_1,
 \qquad |z_0-\rho_1|=a<D.
\]
选择
\begin{equation}
 0<h<\frac12\min\{D-a,\ \sigma-D\}.
 \label{eq:win-shift-margin}
\end{equation}
这是可能的，因为 $D<\sigma$。闭盘
$\overline B(z_0,D+h)$ 仍在 $\mathcal H$ 内，所以
\[
 P=Z\cap\overline B(z_0,D+h)
\]
是有限集。

当 $|t-\operatorname{Im}\rho_1|<h$，令 $z=\sigma+it$，则
\[
 |z-\rho_1|<h+a<D.
\]
此外，$B(z,D)$ 中的每个极点都属于 $P$，因为
\[
 |w-z_0|\le |w-z|+|z-z_0|<D+h.
\]
这样，一段小范围内所有可能影响最近距离的极点，都包含在同一个有限集合 $P$ 中。

对 $P$ 中两个不同点 $p,q$，两者到 $z=\sigma+it$ 的距离平方之差为
\begin{align*}
 &|\sigma+it-p|^2-|\sigma+it-q|^2\\
 &\quad=(\sigma-\operatorname{Re}p)^2
       -(\sigma-\operatorname{Re}q)^2
       +(\operatorname{Im}p)^2-(\operatorname{Im}q)^2
       -2t(\operatorname{Im}p-\operatorname{Im}q).
\end{align*}
如果两点虚部不同，至多一个 $t$ 会使此差为零。
如果两点虚部相同，则实部不同；由两点都严格位于 $\sigma$ 左侧可知
$(\sigma-\operatorname{Re}p)^2\ne
(\sigma-\operatorname{Re}q)^2$，所以没有这样的 $t$。

因此只需从开区间
$(\operatorname{Im}\rho_1-h,\operatorname{Im}\rho_1+h)$
中避开有限多个高度，就能使 $P$ 中各点到 $z$ 的距离互不相同。固定这样的 $t$ 与 $z$，取其中唯一的最近极点 $\rho$。它满足
\[
 0<d:=|z-\rho|<D.
\]
剩余各点的距离都严格大于 $d$；利用有限性，可以选到半径 $R_1$，使
\begin{equation}
 d<R_1<D,\qquad B(z,R_1)\cap Z=\{\rho\}.
 \label{eq:win-isolated-disc}
\end{equation}
若 $P$ 中只有 $\rho$，任取 $d<R_1<D$ 即可。又因
$R_1<D\le\sigma-\theta$，整个圆盘 $B(z,R_1)$ 都在 $\mathcal H$ 内。

到此已固定两个圆盘：$B(\sigma,D)$ 无极点，
$B(z,R_1)$ 只含一个单极点 $\rho$。所有选择均只依赖 $\Phi$、$\theta$ 及上述极点，与随后选取的好尺度 $U$ 无关。

\section{唯一最近极点给出导数下界}

选择一个共同的 Cauchy 半径
\begin{equation}
 d<r<R_1<D\le\sigma-\theta.
 \label{eq:win-cauchy-radius}
\end{equation}
在 $B(z,R_1)$ 内，减去 $\rho$ 处的主部，得到全纯函数
\[
 H(w)=\Phi(w)-\frac{c_\rho}{w-\rho}.
\]
这里在 $w=\rho$ 处取可去延拓；极点是单极点，保证只需减去这一项。

记
\[
 M_1=\max_{|w-z|=r}|H(w)|,\qquad
 M_0=\max_{|w-\sigma|=r}|\Phi(w)|.
\]
两者有限：第一个圆周位于 $H$ 的全纯圆盘内部，第二个圆周严格位于无极点圆盘 $B(\sigma,D)$ 内。Cauchy 导数估计给出，对每个 $k\ge0$，
\begin{align}
 \left|A_k(z)-\frac{c_\rho}{(z-\rho)^{k+1}}\right|
       &\le\frac{M_1}{r^k},\label{eq:win-pole-cauchy}\\
 |A_k(\sigma)|&\le\frac{M_0}{r^k}.
 \label{eq:win-free-cauchy}
\end{align}
第一式还使用了
\[
 \frac{(-1)^k}{k!}
 \left(\frac{c_\rho}{w-\rho}\right)^{(k)}_{\!w=z}
   =\frac{c_\rho}{(z-\rho)^{k+1}}.
\]

这两个估计把极点的贡献与其余全纯部分分开。极点项按 $d^{-k}$ 增长，而其余部分只需用较慢的 $r^{-k}$ 控制；严格距离差 $d<r$ 保证前者最终占优。

由 \eqref{eq:win-pole-cauchy} 和反三角不等式，
\[
 |A_k(z)|\ge\frac{|c_\rho|}{d^{k+1}}-\frac{M_1}{r^k}.
\]
另一方面，
\[
 \operatorname{Re}A_k(\sigma)
       \ge-|A_k(\sigma)|\ge-\frac{M_0}{r^k}.
\]
令
\begin{equation}
 b=\frac{|c_\rho|}{d}>0,\qquad B=M_0+M_1\ge0.
 \label{eq:win-lower-constants}
\end{equation}
于是得到与 \eqref{eq:win-master-upper} 左边完全相同的量的下界：
\begin{equation}
 |A_k(z)|+\operatorname{Re}A_k(\sigma)
       \ge\frac b{d^k}-\frac B{r^k}.
 \label{eq:win-master-lower}
\end{equation}
非零留数保证 $b$ 严格为正。这里无须要求复数留数本身是正实数。

\section{任意远处的好尺度产生矛盾}

现在才让好尺度 $U$ 向无穷远处移动。对每个这样的 $U\ge1$，取
$k=\lceil420U\rceil$。将
\eqref{eq:win-master-upper} 与
\eqref{eq:win-master-lower} 合并，有
\[
 \frac b{d^k}
 \le\frac B{r^k}+\frac{U+C_2}{(\sigma-\theta)^k}.
\]
两边乘以 $d^k$，并记
\[
 q=\frac dr\in(0,1),
\]
由于 $r\le\sigma-\theta$，可得
\begin{align}
 b
 &\le Bq^k+(U+C_2)
          \left(\frac d{\sigma-\theta}\right)^k\notag\\
 &\le(B+U+C_2)q^k
 \le(B+U+C_2)q^U.
 \label{eq:win-final-contradiction}
\end{align}
最后一步使用 $k\ge U$ 与 $0<q<1$。

所有常数 $b,B,C_2,q$ 都已固定，其中 $b>0$。
写 $q^U=e^{-\alpha U}$，$\alpha=-\log q>0$，即可看出
\[
 (B+U+C_2)q^U
       =(B+C_2+U)e^{-\alpha U}\longrightarrow0.
\]
指数衰减超过线性因子的增长。因此存在阈值 $U_*$，使全部 $U\ge U_*$ 的右边都小于 $b$。窗口假设保证仍能选到一个这样的好尺度，与
\eqref{eq:win-final-contradiction} 矛盾。

故不存在 $\rho_0\in Z$ 满足
$\operatorname{Re}\rho_0>\theta$。由于 $\mathcal H$ 中的 $\zeta$ 零点恰好成为 $\Phi$ 的极点，定理
\ref{thm:window-zeta} 得证。\hfill$\square$

这项证明只在选中的窗口内使用算术上界；窗口外的贡献由全局粗界和核尾估计负责。正实轴无极点的性质，则保证可以先找到无极点的实中心圆盘，再通过竖移圆心分离一个非实极点。两项条件各自承担了明确的步骤。

\section{两个反例说明条件不能随意删去}
取 $0<\theta<\beta<1$。若只假定一个实函数有单侧上界，而不检查实极点，那么
\[
g(u)=-e^{\beta u}\le e^{\theta u}
\]
甚至在全部 $u\ge0$ 都成立。但它的拉普拉斯变换为
\[
\Phi(s)=-\frac1{s-\beta},
\]
在 $\beta>\theta$ 有一个留数为负的实极点。因此，“单侧界必然排除所有极点”是错误的。本稿先证明 $\zeta$ 情形没有正实极点，才可以使用后面的圆盘选择。

另取 $g(u)=e^{\beta u}\cos u$，其拉普拉斯变换为
\[
\Phi(s)=\frac{s-\beta}{(s-\beta)^2+1},
\]
在 $\beta+i$ 和 $\beta-i$ 有极点。然而，对
$u=(2n+1)\pi$，有 $g(u)<0\le e^{\theta u}$，这样的好点能出现得任意远。甚至在每段余弦为负的固定长度区间上，上界都成立。
这仍不能排除极点，因为它没有提供本证明需要的、长度随 $U$ 增长的完整窗口 $[U,3U/2]$。

这两个例子说明，我们不能把证明中精确的窗口、实极点条件和导数组合省略后，仅凭“函数有上界”就宣布得到无零点结论。
